\textbf{Tumor-type differences in validated chemicals before and after covariate adjustment.} Each of the 30 validated type-differential chemicals was refit under three nested models: Model 1, unadjusted; Model 2, adjusted for collection year; Model 3, adjusted for collection year, age, and sex. Quantitative features were compared by analysis of covariance with the tumor-type term assessed by an $F$ test, and effect size is partial $\eta^{2}$; qualitative (detection) features were compared by logistic-regression likelihood-ratio test, and effect size is McFadden's pseudo-$R^{2}$. The two effect-size metrics are on different scales and are therefore comparable within a mode but not across modes. Because Fisher's exact test cannot accommodate covariates, likelihood-ratio tests are used throughout this table so that unadjusted and adjusted models are directly comparable; unadjusted values for qualitative features therefore differ slightly from the exact tests reported in Table 3, although every chemical is significant under both. 22 of 30 chemicals remained significant under Model 2, with effect sizes essentially unchanged; the chemicals that lost significance were marginal before adjustment and are not among the findings emphasized in the manuscript. In a sensitivity analysis adjusting instead for the categorical collection intervals of Table 1, 14 of 30 remained significant; because that parameterization discards within-interval variation it is reported here for completeness rather than as a primary specification, and is not tabulated separately. \textit{$\dagger$ Level 2 identification; $\ddagger$ detected in 10 or fewer of the 60 samples (minimum 4), for which covariate-adjusted estimates should be interpreted with caution given the limited number of events per model parameter.}
